\chapter{Tasks, Rewards and Learning System}
\label{ch:learning}

The physical model defines what the rocket can do; the task definition determines what it is
asked to learn. In the prototype, those two concerns were mixed inside \texttt{FalconAgent}.
That was manageable for one landing experiment, but adding hover and target tracking created
conflicting spawn rules, rewards and termination conditions. The learning system was therefore
reorganized around task ownership before the final policies were trained.

\section{Making each task own its objective}

Every scenario now owns a versioned \texttt{ScenarioObjectiveConfig}. It contains the reward
magnitudes, shaping scales and termination parameters that give numerical meaning to that task.
The runtime applies the sign of costs, so values exposed in the UI remain nonnegative magnitudes;
a value of zero disables a term.

The distinction between continuous and one-off rewards became important after early policies
learned to wait for repeated payment. Continuous shaping values are defined as rates and
multiplied by the physics timestep. Contact events and terminal outcomes are emitted once. The
evaluation sequence is also fixed: measure the physical state, build the task context, calculate
shaping, add event and terminal terms, record the breakdown, and only then end the episode. A
terminal event can no longer disappear because another component terminated the episode first.

This organization did not make reward design automatic. It made each assumption visible and
allowed a reward change to become a versioned experimental decision rather than a hidden code
edit.

\section{Putting unlike physical quantities on a common scale}

Position error, speed and attitude are expressed in different units and ranges. Combining their
raw values made coefficients difficult to interpret, so many shaping terms use exponential
closeness
\begin{equation}
  E(x;s)=\exp\!\left(-\frac{\lvert x\rvert}{\lvert s\rvert}\right),
  \label{eq:exp01}
\end{equation}
where the user-visible scale $s$ describes how quickly the value decreases as error $x$ grows.
A general continuous reward rate can then be written
\begin{equation}
  \dot r = \sum_i w_i^+ E(e_i;s_i)-\sum_j w_j^- C_j,
  \qquad r_t=\dot r\,\Delta t+r_{\mathrm{event}}+r_{\mathrm{terminal}}.
  \label{eq:reward-general}
\end{equation}
The $w$ values are configurable magnitudes and each contribution can be logged separately. A
large total return can therefore be traced back to the behaviour that generated it instead of
being treated as proof that the policy is good.

\section{Fixed hover: learning to remain still}

Fixed hover isolates the most basic powered-flight question: can the policy keep the vehicle
near one point without touching the ground or flying away? The rocket starts near
\SI{50}{m} and is asked to settle near a target at \SI{30}{m}. Its reward combines altitude and
planar proximity, uprightness, low translational speed and low angular rate, with optional costs
for motion, effort and time.

This apparently simple task exposed one of the earliest reward loopholes. When similar proximity
terms were used for landing, hovering just above the surface accumulated more reward than ending
the episode with a touchdown. The policy was not refusing to land; the numerical objective had
made waiting more valuable. That behaviour led to the separation of rates, events and terminal
outcomes used by every later task.

A standard hover evaluation lasts \SI{60}{s}. The evaluator reports duration, position,
velocity, attitude and fuel but has no predeclared binary success threshold. The completed
\texttt{HoverFixed} evaluation is therefore reported through its physical distributions and
termination outcomes, not converted into a success percentage after the fact.

\section{Moving-target hover: separating travel from settling}

Once a policy could remain near one point, target tracking added a transition between stable
states. An evaluation episode begins with one hover target and relocates it twice after capture.
Each target must be reached and held within \SI{35}{s}; the curriculum increases the relocation
radius and tightens the capture limits.

A single proximity reward was not enough for both parts of this behaviour. Far from the target,
the useful action is to close distance at a sensible speed. Near it, the same speed becomes a
liability and the policy should settle. The task therefore has an approach phase, which rewards
closure, direction and useful speed, and a settle phase, which rewards centre proximity and calm
horizontal, vertical and rotational motion. Moving away, loitering and excessive speed can be
penalized. Capture is a timed event defined by bounded position and motion, not by crossing a
distance threshold for one frame.

This two-phase design became a useful bridge to landing. Both tasks require the policy to travel
while far away and change priorities as the terminal region approaches.

\section{Physical leg landing: rewarding the whole mission}

Landing adds the final elements missing from tracking: a descent profile, irreversible contact
and a requirement to remain supported afterwards. A height-dependent reference speed is based on
\begin{equation}
  v_{z,\mathrm{goal}}(h)=-k_v\sqrt{2g\max(h,0)},
  \label{eq:descent-profile}
\end{equation}
where curriculum-dependent fraction $k_v$ controls how aggressively the vehicle should descend.
In flight, the objective can use horizontal closure, profile error, uprightness, near-target
speed, angular rate, control effort and time. At the surface, it distinguishes first contact,
hard impact, rebound, support by four feet and stable touchdown.

Runs L2--L11 progressively changed the relationship between those terms. The decisive rule in
the final objective is that mission efficiency is calculated only after a successful landing.
Fuel use, restarts and additional engine ignitions can then distinguish two successful solutions,
but a failing policy cannot improve its outcome by switching everything off and giving up. The
behaviours that led to this rule are followed chronologically in
Section~\ref{sec:landing-iterations}.

Safety after contact is not delegated entirely to the learned policy. Any external impact locks
propulsion, off-pad impact ends immediately, and pad contact remains active long enough to
measure rebound, supported feet and stable hold. Reward is thus a signal used to discover the
behaviour, while evaluation success is a separate conjunction of physical conditions.

\section{Keeping the learning algorithm constant}

The tasks and simulator changed repeatedly, so the principal experiments kept one PPO
configuration instead of adding hyperparameter search as another source of variation.
Table~\ref{tab:ppo} lists the configuration shared by the completed tracking and L11 landing
runs. It should be read as an implementation choice, not as evidence that these values are
optimal.

\begin{table}[htbp]
  \centering
  \caption{PPO configuration for \texttt{HoverTrackSimple3} and \texttt{L11}.}
  \label{tab:ppo}
  \begin{tabular}{ll@{\qquad}ll}
    \toprule
    Parameter & Value & Parameter & Value \\
    \midrule
    Batch size & 4096 & Buffer size & 40960 \\
    Learning rate & $2\times10^{-4}$, linear & PPO clip $\epsilon$ & 0.2 \\
    Entropy $\beta$ & 0.01 & GAE $\lambda$ & 0.98 \\
    Epochs/update & 6 & Discount $\gamma$ & 0.9995 \\
    Hidden units & 512 & Hidden layers & 3 \\
    Time horizon & 1024 & Maximum steps & 30M tracking; 100M landing \\
    Checkpoint interval & 500,000 & Trainer seed & 1 \\
    \bottomrule
  \end{tabular}
\end{table}

The recorded environment used Python 3.10.12, ML-Agents 1.1.0, PyTorch 2.1.0 with CUDA 12.1,
and an NVIDIA GeForce RTX 4070 SUPER. The Unity project used Unity 6000.4.1f1 and ML-Agents
package 4.0.2. Neither principal run used fins, RCS or wind. Tracking used one engine channel,
30 observations and four continuous actions; L11 used three independent engines, 52
observations and 12 continuous actions. With the learning method fixed, the remaining changes
can be interpreted through the task, actuator and simulator revisions described next.
